\documentclass[11pt,a4paper]{article}
\usepackage[T1]{fontenc}
\usepackage[utf8]{inputenc}
\usepackage{lmodern,amsmath,amssymb,amsthm,mathtools}
\usepackage[margin=30mm]{geometry}
\usepackage{enumitem,microtype}
\usepackage[colorlinks=true,linkcolor=blue,citecolor=blue,urlcolor=blue]{hyperref}
\setlist[enumerate]{itemsep=3pt,topsep=5pt}
\setlist[itemize]{itemsep=3pt,topsep=5pt}
\newtheorem{theorem}{Theorem}[section]
\newtheorem{lemma}[theorem]{Lemma}
\newtheorem{proposition}[theorem]{Proposition}
\newtheorem{corollary}[theorem]{Corollary}
\theoremstyle{definition}
\newtheorem{definition}[theorem]{Definition}
\newtheorem{example}[theorem]{Example}
\newtheorem{remark}[theorem]{Remark}
\DeclareMathOperator{\Ty}{Ty}
\DeclareMathOperator{\lv}{lv}
\DeclareMathOperator{\rk}{rk}
\DeclareMathOperator{\FV}{FV}
\DeclareMathOperator{\Sub}{Sub}
\DeclareMathOperator{\Fix}{Fix}
\newcommand{\BH}{\mathrm{BH}}
\newcommand{\ID}{\mathrm{ID}_1}
\newcommand{\N}{\mathbb N}
\newcommand{\W}[2]{#1\succeq_\star #2}
\newcommand{\St}[2]{#1\succ_\star #2}
\newcommand{\EW}[2]{#1\succeq_e #2}
\newcommand{\ES}[2]{#1\succ_e #2}
\newcommand{\cw}{\succeq_c}
\newcommand{\cs}{\succ_c}
\newcommand{\mk}{\mathsf m}
\newcommand{\headmark}{\mathsf m^\circ}
\newcommand{\ap}[2]{#1\langle #2\rangle}
\newcommand{\sharpens}{\mathrel{\#}}
\newcommand{\DM}{\mathsf{DM}}
\newcommand{\Acc}{\mathsf{Acc}}
\newcommand{\It}{\mathsf I}
\newcommand{\tr}{\operatorname{tr}}
\title{A Decidable Path Ordering for G\"odel's T\\
with Bachmann--Howard Descending Height}
\author{}
\date{17 November 2026\\Working manuscript}
\begin{document}
\maketitle
\begin{abstract}
We define a context- and substitution-stable termination ordering for the
literal lambda formulation of G\"odel's T. The comparison uses persistent
marking, simultaneous application along the original lambda spine, and
inclusion of the finite type support of the target in that of the source.
Both weak and strict comparison are decidable, every original beta, case
and recursor reduction is strict, and a disjoint ranked ground sublanguage
has exactly multiset recursive path ordering. The descending height of the
same relation is the Bachmann--Howard ordinal. A finite observation model
proves decidability without bounding generated terms. A numbered simulation
with strictly decreasing finite marks supplies an individual upper bound
through a positive arithmetic accessibility proof in $\ID$. Explicit native
terms realize finite-type ordinal iteration and a cofinal sequence of lower
ranks. The proof distinguishes the structural ordering from comparison of
the numerical ordinal assignments associated with G\"odel's T.
\end{abstract}
\tableofcontents

\section{Introduction and principal theorem}
Termination of G\"odel's T is compatible with several substantially
different ways of measuring terms. Wilken and Weiermann ask for a
higher-order extension of recursive path ordering and anticipate
Bachmann--Howard strength; they also ask about effective comparison of
their ordinal assignments \cite[Introduction]{ww}. This article addresses
the structural ordering question. Its objects are literal typed terms,
including open terms with functional variables.

The ordering is inspired by the marked recursive comparisons and numbered
computability argument developed by Kop \cite[Chapter~5]{kop}. We specify
the mutual rules completely and prove their properties for the relation
defined here. In particular, a mark persists when the strict rule is used
at an already-marked source. This local choice matters in composition:
selected arguments may themselves contain copies of a marked owner.
The well-foundedness proof pays each reuse by a strict decrease of a
finite numerical label.

There are three independent proof obligations. A recursive inference
system need not give a terminating comparison algorithm. A well-founded
relation need not be stable under contexts and substitutions. Finally,
an ordinal lower bound is relevant only when its terms and comparisons
belong to the relation whose upper bound has been proved. The construction
below meets these obligations for one explicitly defined pair.

Fix a finite ranked ground signature $\mathcal F$ with a strict total
precedence. Extend literal T by its actual ranked constructors. All statuses
are multiset statuses. The mixed precedence and literal support $\Ty$
are defined in Section~\ref{sec:syntax}. Sections~\ref{sec:comparison}
and~\ref{sec:stability} define $\cw,\cs$.

\begin{theorem}[Principal theorem]\label{thm:main}
The confined pair $(\cw,\cs)$ has the following properties.
\begin{enumerate}
\item Both comparisons are decidable on all finite ordinary typed inputs.
\item The weak part is reflexive and transitive. The strict part is
transitive and well founded, is included in the weak part, and satisfies
$\cw\,;\,\cs\subseteq\cs$ and $\cs\,;\,\cw\subseteq\cs$.
\item Both parts are preserved by every original well-typed context and
every type-preserving capture-avoiding substitution by ordinary terms.
\item Every original beta, case and recursor reduction is strict.
\item On the first-order $\mathcal F$-algebra, $\cs$ is exactly multiset
recursive path ordering for the chosen precedence.
\item Every individual descending strict rank is below $\BH$, and
\[
 \sup_s\bigl(\rk_{\cs}(s)+1\bigr)=\BH.
\]
An effective family of closed native T terms of outer type $o\to o$
already has cofinal ranks; $\mathcal F$ is unnecessary for this lower bound.
\end{enumerate}
\end{theorem}

Here descending rank means
$\rk_R(s)=\sup_{sRt}(\rk_R(t)+1)$ for a well-founded relation $R$.
The theorem concerns the height of a partial ordering. It does not assert
that the terms form a linear well-order or that every linear extension
has the same order type. No efficiency estimate for the decision procedure
is part of the theorem.

\section{Typed terms and literal type support}\label{sec:syntax}

We use the lambda formulation of G\"odel's system $T$ with the argument
order of Wilken and Weiermann~\cite[Section~1]{ww}.  Its simple types are
generated by
\[
  \sigma,\tau ::= o\mid \sigma\to\tau,
  \qquad
  \lv(o)=0,\quad
  \lv(\sigma\to\tau)=\max\{\lv(\sigma)+1,\lv(\tau)\}.
\]
Arrows associate to the right.  For each type there is a countably infinite,
effectively named collection of variables of that type.  An ordinary term
is a well-typed finite expression formed from variables, abstraction,
application, and the constants
\[
\begin{aligned}
 0&:o,& S&:o\to o,\\
 D^\tau&:o\to\tau\to\tau\to\tau,&
 R^\tau&:o\to(o\to\tau\to\tau)\to\tau\to\tau.
\end{aligned}
\]
The typing rules for variables and constants use their declared types;
from $A:\sigma\to\tau$ and $B:\sigma$ one forms $AB:\tau$, and from
$A:\tau$ with $x:\sigma$ bound in its context one forms
$\lambda x:\sigma.A:\sigma\to\tau$.  Application associates to the left.
Terms are identified up to renaming bound variables.  This identification
does not include beta reduction, eta conversion, or any rearrangement of
application nodes.  Capture-avoiding substitution can be implemented by
fresh renaming or, equivalently, by a typed de Bruijn representation.

The primitive reductions are
\begin{align*}
 (\lambda x:\sigma.A)B&\longrightarrow A[x:=B],\\
 D^\tau 0\,A\,B&\longrightarrow A,
 &D^\tau(Sn)\,A\,B&\longrightarrow B,\\
 R^\tau 0\,A\,B&\longrightarrow B,
 &R^\tau(Sn)\,A\,B&\longrightarrow
      A\,n\,(R^\tau n\,A\,B).
\end{align*}
Their well-typed instances are closed under every well-typed one-hole
context, including contexts below abstractions.  In particular the first
argument of $R^\tau$ is its natural-number test, the second is its step
functional, and the third is its initial value.

\subsection{A ranked first-order sublanguage}

To state an exact recursive-path-order restriction, fix independently a
finite signature $\mathcal F$.  Each $f\in\mathcal F$ has a fixed declaration
$[o,\ldots,o]\to o$ and a fixed number of arguments.  We add the actual
ranked term former $f(A_1,\ldots,A_r):o$.  The names in $\mathcal F$ are
disjoint from the native constants and from the application symbols.
The first-order $\mathcal F$-terms are generated by these ranked formers
and variables of type $o$; they need not be closed.

For the higher-order comparison, ordinary application is represented
literally by the typed binary symbol
\[
 @_{\sigma,\tau}:[\sigma\to\tau,\sigma]\to\tau,
 \qquad
 AB=@_{\sigma,\tau}(A,B).
\]
Every native constant is a nullary head with its declared simple type.
For example $D^o0AB$ consists of a nullary $D^o$ node and three actual
application nodes.  It is not a ternary ranked case symbol.  The distinction
keeps both the native reduction rules and the first-order restriction on
their stated syntactic carriers.

All head statuses in this paper are multiset statuses.  The heads
$@_{\sigma,\tau}$ form one precedence-equivalence class, all native
constant instances form a second class strictly above it, and the heads
of $\mathcal F$ lie above both classes in a chosen strict total precedence.
This is a well-founded precedence with equivalence classes; no comparison
between native type instances is hidden in the precedence.

\subsection{The finite support used for confinement}

For a type $\sigma$, let $\operatorname{Sub}(\sigma)$ be the finite set
of its type subexpressions, including $\sigma$.  Thus
\[
 \operatorname{Sub}(o)=\{o\},\qquad
 \operatorname{Sub}(\sigma\to\tau)
 =\{\sigma\to\tau\}\cup\operatorname{Sub}(\sigma)
                         \cup\operatorname{Sub}(\tau).
\]

\begin{definition}[Literal type support]\label{def:typesupport}
The support $\Ty(A)$ of an ordinary term is the union of
$\operatorname{Sub}(\sigma)$ over the result type $\sigma$ of every
node of $A$ and the declared type of every bound variable of $A$.
Unused entries of an ambient typing context are not included.
\end{definition}

The support is finite and closed under type subexpressions.  At a ranked
ground node, the only declaration types are its actual ground argument
and result types.  In particular its absent curried type is not added to
the support.  At a native nullary head its complete original functional
type is an actual node type and is therefore included.  Auxiliary marked
and numbered terms will retain that original head declaration even when
their current output is retyped; their finite-universe condition will
check both pieces of information.

Write $\operatorname{FV}(A)$ for the finite set of free typed variable
names of $A$.  A finite type-preserving substitution $\theta$ assigns
ordinary terms of the same type to finitely many names and acts as the
identity at all other names.  Its images may contain new free names.

\begin{lemma}[Exact support under substitution]
\label{lem:typesupport-substitution}
For every ordinary term $A$ and finite type-preserving substitution $\theta$,
\[
 \Ty(A\theta)=\Ty(A)\cup
       \bigcup_{x\in\operatorname{FV}(A)}\Ty(\theta x).
\]
\end{lemma}
\begin{proof}
Induct on the literal term.  At a variable $x:\sigma$, the result type of
$\theta x$ is still $\sigma$, so
$\operatorname{Sub}(\sigma)\subseteq\Ty(\theta x)$.  This gives the
displayed equality, including the identity case.  A constant keeps its
type and has no free names.  At application and at a ranked node, the
node's result and all argument types are unchanged, and the induction
hypotheses combine by finite union.  At an abstraction, rename the binder
fresh for the finitely many relevant substitution images and mask its
name.  Its declaration and the abstraction type are retained, while
the induction hypothesis for the body contributes precisely the free
names of the abstraction.  Bound renaming changes no type support.
\end{proof}

\begin{corollary}[Support confinement is stable]
\label{cor:typesupport-stability}
Suppose $A,B$ have the same outer type,
$\Ty(B)\subseteq\Ty(A)$, and
$\operatorname{FV}(B)\subseteq\operatorname{FV}(A)$.
Then $\Ty(B\theta)\subseteq\Ty(A\theta)$ for every such $\theta$.
Moreover $\Ty(C[B])\subseteq\Ty(C[A])$ for every well-typed one-hole
context $C$ of the appropriate hole type.
\end{corollary}
\begin{proof}
The substitution assertion follows by inclusion in both unions of
Lemma~\ref{lem:typesupport-substitution}.  A fixed context contributes
the same finite set of node and binder types to both plugged terms;
the two hole types agree.  Taking the union with those common types
preserves support inclusion.
\end{proof}

This support condition permits one finite type universe to control every
descendant of a fixed source.  It does not bound the size of descendants,
the length of generated argument lists, or the types that may occur in
other, unrelated input terms.

\section{Persistent marking and recursive comparison}\label{sec:comparison}

The ordinary syntax is enlarged only for the purpose of defining comparison.
An ordinary head $f$ retains its declared base arguments
$\vec s=(s_1,\ldots,s_r)$. A marked head has the form
\[
 P_\tau=f^\star_\tau(\vec s;\vec e).
\]
Its base arguments still have their declared types. Its finite list of
extras $\vec e$ has arbitrary recorded types, and its output annotation
$\tau$ can differ from the original result type of $f$. A retyping $P_\rho$
changes only that output annotation. Variables and lambdas are retained,
and their subterms may contain marks. An \emph{ordinary} term has no marked
node anywhere. Every judgment below compares terms of equal outer type.

\begin{definition}[Persistent partial marking]
Define $\mk(x)$ to be undefined,
$\mk(\lambda x.s)=\lambda x.\mk(s)$ when the body mark exists, and
\[
 \mk(f(\vec s))=f^\star(\vec s;\varnothing),\qquad
 \mk(f^\star_\tau(\vec s;\vec e))=f^\star_\tau(\vec s;\vec e).
\]
Write $\headmark(s)$ for total top marking, which leaves a lambda spine
ending in a variable unchanged and otherwise agrees with $\mk$.
Thus $(\headmark)^2=\headmark$.
\end{definition}

\begin{definition}[Original-spine application]\label{def:application}
For $q:\rho_1\to\cdots\to\rho_n\to\tau$, the partial expression
$\ap q{a_1,\ldots,a_n}$ is defined as follows. An empty tuple returns $q$.
Consume its original initial lambdas, assigning the supplied arguments
simultaneously to their bound variables. If all arguments are consumed,
return the substituted residual term. If arguments remain, inspect the
original residual head. An ordinary or marked functional head is retained,
marked if necessary, its old arguments are substituted, and all remaining
arguments are appended to its extras at once. An original variable head
blocks the operation, even when its substituted value is functional.
\end{definition}

This last qualification is essential. With $\theta=o\to o$, a functional
head $h:\theta$, and $y:o$, the full call
$\ap{(\lambda x:\theta.x)}{h,y}$ is undefined, whereas the two successive
unary calls return $h^\star(;y)$. A successful full call does commute
with capture-avoiding substitution. This follows by following its original
lambda spine: the only branch that could inspect a substituted variable
is excluded by the success assumption.

\begin{definition}[Occurrence multiset extension]\label{def:multiset}
For a weak relation $W$ and a strict relation $S$, write
$A>_{\mathrm{mul}}^{W,S}B$ when the source occurrences split into
$A_0\uplus A_1$ with $A_1\ne\varnothing$, the target occurrences split
into $B_0\uplus B_1$, $A_0$ and $B_0$ are bijectively matched by
equal-typed $W$ judgments, and every occurrence of $B_1$ is $S$-below
some occurrence of $A_1$. A source occurrence may supply several new
occurrences; it may also supply none. Equal terms remain distinct
occurrences.
\end{definition}

Define $\succeq_\star,\succ_\star$ simultaneously as the least relations
closed under the following clauses. Every derivation is finite. In the
clauses concerning a marked owner, $P=f^\star(\vec s;\vec e)$; all copies
are retyped to their indicated types. Equivalent heads are denoted by
$f\approx g$, and strict head precedence by $f>g$.
\begin{enumerate}[label=\textsc{\arabic*.}]
\item[\textsc{Var}] $\W{x}{x}$.
\item[\textsc{Abs}] From $\W{s}{t}$ derive
$\W{\lambda x.s}{\lambda x.t}$, opening both binders with the same fresh name.
\item[\textsc{Fun}] $\W{f(\vec s)}{g(\vec t)}$ if the heads are equivalent,
their result types agree, and a permutation of $\vec t$ is pointwise
weakly matched with $\vec s$.
\item[\textsc{Put}] From $\W{f^\star(\vec s;\varnothing)}{t}$ derive
$\W{f(\vec s)}{t}$.
\item[\textsc{Select}] If $q$ is a base or extra argument of $P$, has type
$\rho_1\to\cdots\to\rho_n\to\tau$, and
\[
 \ap q{P_{\rho_1},\ldots,P_{\rho_n}}=r,\qquad \W r t,
\]
then $\W{P_\tau}{t}$. The tuple must succeed as one full call.
\item[\textsc{FAbs}] If
$\W{f^\star_\tau(\vec s;\vec e,x)}{t}$ after fresh opening at
$x:\sigma$, then
$\W{P_{\sigma\to\tau}}{\lambda x:\sigma.t}$.
\item[\textsc{Copy}] If $f>g$ and $\W{P_{\operatorname{type}(t_i)}}{t_i}$
for every target base argument, then $\W{P_{\operatorname{res}(g)}}{g(\vec t)}$.
\item[\textsc{Stat}] The same conclusion holds if $f\approx g$,
$\vec s>_{\mathrm{mul}}^{\succeq_\star,\succ_\star}\vec t$, and
all the same whole-source weak premises hold.
\end{enumerate}
The strict rule is
\[
 \frac{\mk(s)=s'\qquad \W{s'}{t}}{\St{s}{t}}.
 \tag{Strict}
\]
The base list, including its multiplicities, is used in \textsc{Stat}.
Extras are selectable but are absent from that multiset test.

\begin{remark}[Why persistence is material]
Let $Z<H<K<F$ be ground ranked heads of arities $0,0,1,2$, let $k=K(H)$,
and put
\[
 P=(\lambda x.F(x,Z))k,\qquad T=F(k,Z),\qquad U=F(H,H).
\]
Marking $P$ gives $q=@^\star(\lambda x.F(x,Z),k)$. Selection gives
$\W q H$, and persistence therefore gives $\St q H$. In
$F^\star(q,Z)$ remove both base occurrences and replace them by two
$H$ occurrences, each supplied strictly by $q$. The whole-source premises
to $H$ follow by selection of $q$. Thus \textsc{Stat} and \textsc{Put}
give $\W{F(q,Z)}U$. Selecting the lambda argument of $q$ with one
self-copy gives this very term, so $\St P U$.
This derivation uses an already-marked source in a strict multiset premise.
A definition that makes such a source unmarkable cannot use this derivation.
Our proofs below are for the persistent definition just given.
\end{remark}

\section{Structural laws, confinement and original reductions}
\label{sec:stability}

\begin{lemma}[Targets, inclusion and reflexivity]\label{lem:elementary}
Every target of $\succeq_\star$ or $\succ_\star$ is ordinary.
Strict comparison implies weak comparison. Every ordinary term is
weakly reflexive.
\end{lemma}
\begin{proof}
For the first assertion induct jointly over weak, strict and occurrence
list derivations. The variable target is ordinary. \textsc{Abs} and
\textsc{Fun} use their child targets; \textsc{Put} and \textsc{Select}
retain a premise target. \textsc{Copy} and \textsc{Stat} have an ordinary
target head and weak premises to every target argument. \textsc{FAbs}
binds its ordinary body target. The strict rule inherits its weak target.

If $\mk(s)=s'$ and $\W{s'}t$, peel the common initial lambdas.
At an ordinary head \textsc{Put} yields $\W s t$; at a marked head
$s'=s$. A variable cannot be marked. Lifting by \textsc{Abs} gives
strict inclusion. Finally structural induction on an ordinary term,
using \textsc{Var}, \textsc{Abs}, and identity matching in \textsc{Fun},
proves reflexivity.
\end{proof}

\begin{lemma}[Substitution and contexts]\label{lem:structural-stability}
The pair on ordinary endpoints is preserved by ordinary type-preserving
substitution and by every original well-typed context.
\end{lemma}
\begin{proof}
A successful mark commutes with substitution, including its marked case:
$\mk(s\theta)=\mk(s)\theta$. The equality follows by induction down
the initial lambda spine; substitution changes neither an ordinary head
nor a marked head. Successful full application also commutes with
substitution, by Definition~\ref{def:application}.

Induct simultaneously over weak, strict, weak-list matching and strict
replacement derivations. \textsc{Var} uses ordinary reflexivity of
$\theta x$. Under a binder use the lifted substitution. Head rules
transport their argument witnesses. \textsc{Select} uses the successful
application identity and the identity between substituted self-copies
and self-copies of the substituted owner. \textsc{FAbs} uses fresh
opening. \textsc{Strict} uses mark covariance. Applying a substitution
to a removed list preserves its nonempty set of occurrences even when
different terms acquire identical substituted values.

Weak context closure follows by induction on the context, using
\textsc{Abs} and \textsc{Fun}. For a strict hole under an ordinary head,
mark that head and use \textsc{Stat}: retain every unchanged occurrence
and replace the distinguished source argument strictly. Select the
corresponding source argument to supply each whole-source side premise,
using reflexivity or strict inclusion as appropriate. Under a lambda
lift the successful mark and its weak proof. Literal application is
one of these ordinary heads, so both of its argument positions are covered.
\end{proof}

\begin{lemma}[Free-variable containment]\label{lem:fv}
If $\W s t$ or $\St s t$, then $\FV(t)\subseteq\FV(s)$.
\end{lemma}
\begin{proof}
Use the same mutual derivation induction. Head matching distributes target
variables among source arguments. Copy and Stat have whole-source premises
to every target argument. A Select result uses only variables of the
selected argument and its owner copies, all already in the owner.
Opening adds one fresh variable, and the target abstraction binds it.
Marking changes no free variable. These observations give each rule case.
\end{proof}

Composition requires more than the preceding structural inductions.
Appendix~\ref{app:composition} gives a conservative auxiliary relation
on marked targets and proves its ordinary-final cut. Its proof preserves
the genuine failure of unrestricted associativity of full application.

\begin{theorem}[Ordinary composition]\label{thm:composition}
On ordinary endpoints, weak comparison is transitive, and
\[
 \succeq_\star\,;\,\succ_\star\ \subseteq\ \succ_\star,\qquad
 \succ_\star\,;\,\succeq_\star\ \subseteq\ \succ_\star.
\]
Consequently strict comparison is transitive.
\end{theorem}
\begin{proof}
Embed both weak proofs in the auxiliary relation of
Appendix~\ref{app:composition} and apply Theorem~\ref{thm:cut}.
For strict-then-weak use $\mk(s)$ as the initial source in the same cut.
For weak-then-strict, Lemma~\ref{lem:extended-marking} transports the
first weak proof to the successfully marked endpoints; markability of
the intermediate term reflects to the source. Cut with the weak witness
for the second strict judgment and apply \textsc{Strict}.
Strict inclusion then gives strict transitivity.
\end{proof}

\begin{definition}[Confined pair]\label{def:confined}
For ordinary terms of equal type define
\[
 s\cw t\iff \W s t\ \land\ \Ty(t)\subseteq\Ty(s),\qquad
 s\cs t\iff \St s t\ \land\ \Ty(t)\subseteq\Ty(s).
\]
\end{definition}

\begin{theorem}[Stable pair]\label{thm:stable}
The weak confined relation is reflexive and transitive, its strict part
is included in it and transitive, and both mixed composition laws hold.
Both parts are stable under original typed contexts and ordinary
capture-avoiding substitutions.
\end{theorem}
\begin{proof}
Support inclusions compose. The relational assertions therefore follow
from Lemma~\ref{lem:elementary} and Theorem~\ref{thm:composition}.
For substitution use Lemma~\ref{lem:fv} and
Corollary~\ref{cor:typesupport-stability}; every image type introduced
on the target is also introduced at a source occurrence. A fixed context
adds the same types to both sides. Combine these facts with
Lemma~\ref{lem:structural-stability}. Well-foundedness is proved separately
in Section~\ref{sec:upper}, without assuming this composition theorem.
\end{proof}

\begin{theorem}[Exact ranked restriction]\label{thm:rpo}
On first-order $\mathcal F$-terms, including patterns with ground-type
variables, $\cs$ is exactly multiset recursive path ordering. Its weak
part consists of that strict ordering together with recursive multiset
equality.
\end{theorem}
\begin{proof}
Every type support is $\{o\}$. An unmarked first-order source uses only
\textsc{Var}, \textsc{Fun}, or \textsc{Put}; a marked one uses only
\textsc{Select}, \textsc{Copy}, or \textsc{Stat}. Every selectable argument
has ground type, so its tuple is empty. No extras, lambdas, or marks
inside base arguments are generated. In particular, strict multiset
premises always have ordinary source arguments.

Induct on a weak derivation, interpreting an ordinary source comparison
as weak RPO and a root-marked source comparison as strict RPO.
\textsc{Var} is equality. In \textsc{Fun}, either every matched pair
is recursively equal or at least one is strict, giving a strict multiset
decrease. \textsc{Put} includes strictness into weakness. Selection is
the subterm clause, Copy the strict-precedence clause, and Stat the
multiset clause. The whole-source side premises are precisely the usual
domination requirements.

Conversely use \textsc{Var}/\textsc{Fun} for recursive multiset equality
and induct on an RPO derivation. The subterm clause is empty-tuple
Select. Strict precedence uses Copy, and the same-head clause uses Stat.
Each target argument is reached from its weakly matched or strictly
dominating source occurrence by selection, using strict inclusion.
Thus every side premise is available. The two inductions take place
entirely in the ranked fragment.
\end{proof}

\begin{lemma}[Ordinary template grafting]\label{lem:grafting}
If $A$ is an ordinary template and $\W p q$, then
$\W{A[x:=p]}{A[x:=q]}$, even when $p$ contains marks.
\end{lemma}
\begin{proof}
Induct on the ordinary template. At its distinguished variable use the
given comparison; other variables use Var. Abstraction uses fresh
opening, and every ordinary head uses Fun with identity argument
matching. The target image $q$ is ordinary by
Lemma~\ref{lem:elementary}, so the target remains ordinary.
\end{proof}

\begin{theorem}[Original reduction orientation]\label{thm:orient}
Every beta, $D0$, $DS$, $R0$, and $RS$ reduction, including its context
and substitution instances, strictly decreases under $\cs$.
\end{theorem}
\begin{proof}
For a beta redex $(\lambda x:\sigma.A)B:\tau$, put
$P_\rho=@^\star_\rho(\lambda x.A,B)$. Select its second base argument
to obtain $\W{P_\sigma}B$. Grafting gives
$\W{A[x:=P_\sigma]}{A[x:=B]}$. Selecting the first base argument
with its one self-copy produces exactly the former term. This gives
the required marked weak proof and hence strictness.

For $D^\tau uAB$, the outer second argument is $B$, so Select reaches
$B$. To reach $A$, select the outer first argument with one self-copy;
the resulting marked inner application still has $A$ as its second
base argument, selected with an empty tuple. These comparisons work
for every $u$, hence in particular for the two case rules. The $R0$
rule similarly selects the outer base argument.

For the successor recursor put $H=o\to\tau\to\tau$,
$U_1=R^\tau(Sn)$, $U_2=U_1A$, $L=U_2B$ and $L'=R^\tau nAB$.
Use the marked families
\[
\begin{aligned}
 P_\rho&=@^\star_\rho(U_2,B),\\
 C_{2,\rho}&=@^\star_\rho(U_1,A;P_\tau),\\
 C_{1,\rho}&=@^\star_\rho(R^\tau,Sn;C_{2,H},C_{2,\tau}),\\
 C_{0,\rho}&=(R^\tau)^\star_\rho(;C_{1,o},C_{1,H},C_{1,\tau}).
\end{aligned}
\]
Three Select prefixes carry a weak proof from $C_{0,\tau}$ back to
$P_\tau$, selecting $U_2,U_1,R^\tau$ with respectively one, two and
three simultaneous self-copies at their actual types.
The extra chain from $C_{0,H}$ reaches $A$. The chain from $C_{0,o}$
reaches $Sn$ and then $n$. Here $\St{Sn}n$ is itself outer Select.
Context closure of this strict comparison gives $\St L{L'}$, whose
weak witness is $\W{P_\tau}{L'}$. Prefixing the extra chain therefore
gives $\W{C_{0,\tau}}{L'}$.

Two Copy clauses, using the native head precedence $R^\tau>@$, now
construct first $An$ at type $\tau\to\tau$ and then $AnL'$ at type
$\tau$. Prefixing the three outer selections gives
$\W{P_\tau}{AnL'}$, as required. No native recursor has been treated
as a ranked ternary symbol.

Each contractum uses only node, binder and application-suffix types
already in its redex. Hence every displayed strict comparison is
confined. The stable-pair theorem gives all context and substitution
instances.
\end{proof}

\section{An exact finite comparison procedure}\label{sec:decision}

Selection can insert copies of an owner into one of its arguments. Neither
the size of the generated source nor the length of its extra list has an
evident bound. The procedure below instead bounds the information that a
fixed target can observe. The full tuple operations are part of that
information; replacing them by iterated unary application would be unsound.

Fix ordinary inputs $s,t$. Let $K$ be the subtype closure of their literal
types, let $V$ be their free-name set, and let $B$ be the maximum lambda
nesting depth in $t$. Use $B$ fresh port slots, each with a type in $K$.
A context $\Gamma$ is a finite assignment of types to used slots. Its
extensions add unused slots without changing old types.

For each $\Gamma,\sigma$, let $O_\sigma^\Gamma$ consist of all opened
subterms of $t$ of type $\sigma$ whose externally bound variables are
assigned to appropriately typed ports and whose remaining lambda depth is
at most $B-|\Gamma|$. This is a finite effectively constructed family.
Opening a target lambda consumes one unused port and descends to one of
these observations.

\begin{lemma}[Finite closure]\label{lem:observation-closure}
In a recursive comparison originating at $(s,t)$, every requested type
belongs to $K$, every source head occurs in $s$, every target belongs to
the corresponding observation family, and every free source name belongs
to $V\cup\Gamma$.
\end{lemma}
\begin{proof}
Select requests domain and suffix types of a selected argument; these
are subtypes of its existing type. Its owner copies change only their
output annotation. Copy and Stat request types of target arguments.
FAbs uses the domain and codomain of the observed lambda. Substitution
keeps body-node types. No rule creates a new source head: in particular
the target head introduced by Copy is not a new recursive source.
Only descent through a target lambda introduces a port, so used ports
plus remaining target depth never exceeds $B$. Mark reuse introduces
no node or name. Joint induction over the comparison and list rules
proves the assertions.
\end{proof}

\subsection{Natural finite domains}
At each type $\sigma$ and context $\Gamma$ construct a finite domain
$D_\sigma^\Gamma$. A value contains, for each extension $\Delta$ of
$\Gamma$, an availability bit and weak and strict subsets of
$O_\sigma^\Delta$. If $\sigma$ is an arrow type, it additionally contains
a full tuple table for every nonempty suffix decomposition
$\sigma=\rho_1\to\cdots\to\rho_n\to\tau$:
\[
 \operatorname{App}_\Delta:
 D_{\rho_1}^\Delta\times\cdots\times D_{\rho_n}^\Delta
 \longrightarrow D_\tau^\Delta .
\]
Every argument and result type here is a proper type subexpression, so
this is a recursion on type size. Tables are monotone. Order all data
pointwise, and let weakening $a\mapsto a|_\Theta$ retain precisely the
components indexed by extensions of $\Theta$. Retain only values satisfying
\[
 \operatorname{App}_\Delta(a,\vec v)|_\Theta
 =\operatorname{App}_\Theta(a|_\Theta,\vec v|_\Theta)
 \quad(\Delta\subseteq\Theta).
 \tag{N}
\]
The domains are hereditary: the values in the table entries already belong
to the natural lower-type domains. The bottom value has false availability,
empty observation sets, and constantly bottom tuple tables.

\begin{lemma}\label{lem:finite-lattices}
These domains are effective finite complete lattices. Weakening preserves
joins and meets.
\end{lemma}
\begin{proof}
At base type use finite products of bits and powersets. At an arrow, the
lower-type domains are already available. Enumerate all finite tables and
retain the monotone natural ones. Pointwise joins and meets preserve
monotonicity and (N), since weakening is a coordinate projection and
preserves lower-type joins and meets. The all-bottom and all-top families
satisfy the same equations. This proves both the inductive construction
and the lattice assertions.
\end{proof}

For each possible source head $f$, output type $\tau\in K$, and context
$\Gamma$, use a monotone marked-constructor table
\[
 M_{f,\tau}^\Gamma(\vec b,E)\in D_\tau^\Gamma.
\]
The base tuple has $f$'s original declared types. The extra argument $E$
is a down-set of the finite disjoint union of the typed domains at
$\Gamma$. Require coherence
\[
 M_{f,\tau}^\Gamma(\vec b,E)|_\Delta
 =M_{f,\tau}^\Delta
   \bigl(\vec b|_\Delta,\downarrow\{a|_\Delta:a\in E\}\bigr).
 \tag{C}
\]
The coherent monotone families form another finite complete lattice by
the pointwise argument. Base positions and multiplicities are retained.
Extras use a down-set because the rules select from them but never count
their occurrences or use their order in a status test.

For clarity, permutation or repetition of extras already preserves
derivability syntactically. Induct on a finite derivation: Select finds
the corresponding occurrence, FAbs appends the same fresh variable,
and Copy/Stat inspect only the unchanged base list. The remaining rules
transport child proofs. This observation explains the set representation;
the adequacy proof below also proves that taking downward closure loses
no comparison.

\subsection{Structural interpretation and the operator}
At a constructor family $M$, interpret a finite source term in an
environment of abstract values. An undefined result is interpreted by
the entire bottom value, including availability and tuple tables.
For a marked state $(f,\vec b,E)$ compute its weak observation set by
the following four positive tests, first weakening to the observation's
context.
\begin{enumerate}
\item Select a base value or an element of $E$, choose a compatible
suffix decomposition, and apply its one full tuple table to retyped
$M$-values of the same owner. Inspect the resulting weak target bit.
\item For Copy, test the strict head precedence and the weak bits of
the retyped whole owner against all target arguments.
\item For Stat, test equivalent heads, the finite occurrence multiset
witness on $\vec b$, and the same whole-owner weak bits.
\item For FAbs, open the observed lambda at a fresh port, weaken the
base and extras, and add the port value to the extra down-set.
\end{enumerate}
Set availability to true and set \emph{both} the weak and strict fields
equal to this newly computed weak set. The strict equality records
$\mk(P)=P$; no strict bit is created without a weak rule witness from
the previous approximation. The full tuple table appends the entire
tuple to $E$ and consults $M$ at the resulting output type.

For an ordinary head, weak bits arise from Fun or from its marked
counterpart; strict bits are the weak bits of that counterpart. Nonempty
tuple tables mark its head and append all arguments simultaneously.
A literal variable has its matching weak bits, no strict bits and
undefined nonempty tuple tables.

A lambda has weak and strict bits against lambda observations given by
the corresponding opened body fields. To calculate a tuple table, scan
its \emph{original} lambda spine and record consumed binders as protected.
When arguments remain at a protected-variable head, return bottom.
At a functional head append every remaining argument at once.
An unprotected free variable supplied by the abstract environment uses
its recorded full tuple table. If no argument remains, interpret the
residual term under the updated environment. Bound renaming uses fresh
names outside the finite free-name and port sets; it introduces no new
observation context.

\begin{lemma}[Structural interpretation]\label{lem:structural-interpretation}
Interpretation and this original-spine helper are monotone and commute
with weakening. Interpretation respects capture-avoiding substitution.
For environments interpreting concrete substitutions, the helper
interprets full application of the resulting substituted source,
including its undefined cases. This is not a claim that a previously
undefined application remains undefined under free substitution.
\end{lemma}
\begin{proof}
Induct simultaneously on the source and its original-spine helper.
A variable uses its assigned natural value; a protected-variable call
with remaining arguments is constantly bottom. At a head use
monotonicity and (C). At a lambda, fresh renaming and the extended
environment reduce the claim to its original body. Matching bits
against lambda observations remain independent of tuple behavior.

If a tuple ends within the original lambda spine, the two calculations
are ordinary simultaneous substitution. Otherwise a protected bound
variable blocks both calculations; an unchanged functional head gives
the same mark, retyping and appended tuple. An unprotected free variable
is first replaced by its substitution image, whose full table records
that image's original spine. Thus this case may succeed even when the
unsubstituted variable's call would fail. Naturality is (N) in that
variable case and (C) at a constructor. All uses of function tables are
on the entire remaining tuple; no unconditional currying equation is used.
\end{proof}

The preceding equations define an operator $\mathcal F$ on the finite
constructor lattice.

\begin{lemma}\label{lem:operator}
$\mathcal F$ is a monotone endomap of the coherent natural families.
Its iteration from bottom reaches the least fixed point in finitely
many steps.
\end{lemma}
\begin{proof}
Every rule test is positive in weak and strict bits. A multiset test
asks for a finite occurrence assignment with positive premises, so
adding bits preserves an existing witness. Monotone tuple application,
weakening and downward closure preserve this property. Copying a
computed weak set into the strict field is also monotone.

At an extension $\Delta$, the observation equations consult only
weakened inputs and $M^\Delta$. Tuple outputs respect coherence because
\[
 \downarrow(E\cup\{v_i\})|_\Delta
 =\downarrow\bigl(E|_\Delta\cup\{v_i|_\Delta\}\bigr),
\]
where weakening of a down-set means downward closure of its image.
Both inclusions follow from monotonicity of weakening and the presence
of the original image generators. This proves (N) and (C) for the
new family. Finiteness now gives effective stabilization of the
ascending bottom iteration.
\end{proof}

\subsection{Adequacy}
\begin{theorem}[Exact observation semantics]\label{thm:adequacy}
At the least fixed point, a weak or strict observation bit is true
exactly when its corresponding finite recursive judgment is derivable.
\end{theorem}
\begin{proof}
For soundness relate an abstract value $a$ at $\Gamma$ to a concrete
typed source $q$ with $\FV(q)\subseteq V\cup\Gamma$. At every extension,
all weak and strict bits must be genuine finite proofs from $q$.
Every full tuple table must send related concrete arguments to a value
related to their exact special-application result. Only bottom relates
to an undefined result. By type induction this relation is downward
closed; bottom relates to every permitted defined term.

Prove by stage induction that each constructor approximation preserves
this relation on related base arguments and extra down-sets. At bottom
there are no bits. For a new Select bit, a selected base value has its
given concrete argument. A selected extra value lies below the abstract
value of an actual extra and therefore relates to that same concrete
extra. The full tuple property supplies the literal Select premise.
The head clauses use their displayed concrete premises; in particular
Stat's strict bits are genuine proofs of the present strict relation.
FAbs uses the fresh port and the opened source in the enlarged context.
The tuple output is a concrete marked owner with all arguments appended.
The joint structural induction of
Lemma~\ref{lem:structural-interpretation} supplies interpretation of
ordinary heads, lambdas, and substitution environments.

The new marked strict bit is justified by the same already constructed
weak proof $\W P u$ as its weak bit and the equation $\mk(P)=P$.
One application of Strict therefore supplies it. This reasoning also
lifts through every initial lambda. It uses no composition or
well-foundedness premise. Thus every finite stage, and hence the stable
stage, is sound.

For completeness induct jointly on finite weak, strict and list
derivations. Var, Abs and Fun are the structural equations. Put inserts
the marked weak bit into the ordinary weak field. Copy, Stat and FAbs
are exactly their positive equations. Select uses the exact full tuple
table; every actual extra occurs in its downward closure.
Lemma~\ref{lem:observation-closure} retains every requested observation.
For Strict, its shorter weak premise appears by induction. At the fixed
point the strict field of $q$ is exactly the weak field of $\mk(q)$
when marking succeeds, and is empty otherwise: verify this down the
initial lambda spine, using the ordinary-head equation, the persistent
marked equation, and the empty variable case. This supplies the strict
bit, including when its source already contains marks.
\end{proof}

\begin{theorem}[Decidability]\label{thm:decide}
Both $\succeq_\star$ and $\succ_\star$, and therefore $\cw,\cs$, are
decidable on all ordinary finite inputs.
\end{theorem}
\begin{proof}
Construct the finite domains and coherent constructor lattice, iterate
to its least fixed point, interpret the source at the empty port context,
and inspect the target's weak or strict bit. Theorem~\ref{thm:adequacy}
gives both directions of correctness, and Lemma~\ref{lem:operator} gives
termination. Add the finite literal-support inclusion test for confinement.
The algorithm allows functional free variables and arbitrarily high
finite types. It requires no generated-source size bound, search for an
ordinary intermediate term, or transitive-closure oracle.
\end{proof}

\section{A numerical simulation and the individual upper bound}
\label{sec:upper}

Persistent symbolic marking must not become a numerical identity step.
We simulate each reuse by decreasing a finite label. The argument follows
the computability architecture of numbered higher-order path orderings
\cite[Section~5.3]{kop}, with that additional decrease made explicit.
It proves well-foundedness independently of the composition theorem.

\subsection{One finite arithmetic universe}
Fix a finite subtype-closed set $S$ and the ranked signature $\mathcal F$.
The auxiliary syntax has variables, lambdas, bottom $\bot_\tau$, ordinary
heads, numbered heads $f^k_\tau(\vec b;\vec e)$ for $k\in\N$, and a
separate ordinary auxiliary application $QU$. An ordinary head has label
$\omega$. Auxiliary application is distinct from the literal original
binary head $@_{\sigma,\tau}$.

Every node and binder type lies in $S$. In addition, every head's
\emph{original} result type and all original base-argument types lie in
$S$, even when a numbered head has been retyped. These restrictions give
finitely many possible native and application heads. Ranked heads belong
to the fixed finite signature. An ordinary term is in this universe
whenever its literal support is contained in $S$.

Use de Bruijn bound syntax and globally typed free names $(\tau,i)$,
where $\tau\in S$ and $i\in\N$. All finite contexts and their extensions
are represented in this one arithmetic carrier. Fresh opening chooses
a name outside a finite set. There is no application of a theorem about
one fixed context to unrepresented later context extensions.

Write $P^k_\rho=f^k_\rho(\vec b;\vec e)$. The raw auxiliary reduction
$\leadsto$ has the following root rules, always with the displayed
types and head-window conditions.
\begin{enumerate}
\item An ordinary $f(\vec b)$ reduces to $f^k(\vec b;\varnothing)$
for any finite $k$.
\item $P^{k+1}_\tau\leadsto P^k_\tau$.
\item If a base or extra argument $q$ has type
$\rho_1\to\cdots\to\rho_n\to\tau$, then
\[
 P^{k+1}_\tau\leadsto qP^k_{\rho_1}\cdots P^k_{\rho_n}.
\]
The right side uses ordinary auxiliary application.
\item If $f>g$ and $g$ belongs to the fixed head window, then
$P^{k+1}_{\operatorname{res}(g)}$ reduces to the ordinary $g$-head
whose base arguments are the appropriately retyped copies $P^k$.
\item If $f\approx g$, $g$ lies in the same window, and the base list
of $g(\vec c)$ is a primitive strict replacement of $\vec b$ as defined
below, then $P^{k+1}_{\operatorname{res}(g)}\leadsto g(\vec c)$.
\item $P^{k+1}_{\sigma\to\tau}$ reduces to
$\lambda x:\sigma.f^k_\tau(\vec b;\vec e,x)$ with fresh opening.
\item $(\lambda x.Q)U\leadsto Q[x:=U]$, and every nonbottom $Q:\tau$
reduces to $\bot_\tau$.
\end{enumerate}
Close these rules under lambda bodies and under each base and extra
argument of an ordinary or numbered head. Auxiliary application has
\emph{function-side closure only}: $Q\leadsto Q'$ gives
$QU\leadsto Q'U$. It has no argument-side closure.

For primitive replacement, a \emph{mark decrease} passes through a
common initial lambda spine and either changes an ordinary head's label
$\omega$ to a finite label, or changes a finite leading label $k$ to
$l<k$. It leaves the head, arguments, extras and output type fixed.
Variables, bottoms and auxiliary applications cannot be marked this way.
Every mark decrease is a nonempty raw path. A primitive multiset
replacement retains identical occurrences, removes a nonempty source
part, and inserts finitely many mark decreases of removed occurrences.
Deletion is allowed. In particular label zero cannot be marked again.

Quotient by the congruence that replaces equivalent heads and permutes
their permitted base arguments, preserving types, labels, lambda shape
and extra order. Its classes are finite and effectively enumerable in
the fixed head window. We may use least-code representatives. Let
$R_S$ be the positive-path relation induced by $\leadsto$ on these
classes. Both relations are arithmetic: finite raw paths and finite
equivalence witnesses are numerical certificates.

\subsection{Accessibility in one positive inductive definition}
\begin{lemma}[Arithmetic accessibility]\label{lem:arithmetic-accessibility}
For each fixed $S,\mathcal F$, the theory $\ID$ proves accessibility
of every well-typed auxiliary term under $R_S$.
\end{lemma}
\begin{proof}
Use two positive arithmetic operators, with typedness included:
\[
 I(Q)\ \longleftrightarrow\
   \forall Q'\,(Q R_S Q'\Rightarrow I(Q')),
 \qquad
 J(M)\ \longleftrightarrow\
   \forall M'\,(M\DM(R_S)M'\Rightarrow J(M')).
 \tag{A}
\]
Here $\DM(R_S)$ is ordinary strict finite multiset replacement, retaining
identical occurrences and replacing removed ones by $R_S$ descendants.
Neither clause refers to the other's inductive predicate. Disjoint
numerical tags combine them into one positive arithmetic definition.
The $I$-stages of that definition equal the standalone accessibility
stages, since its clause never consults the $J$ tag.

For each of the finitely many types in $S$ define, by external recursion
on type,
\[
 C_o(Q)=I(Q),\qquad
 C_{\sigma\to\tau}(Q)\iff
       \forall U\,(C_\sigma(U)\Rightarrow C_\tau(QU)).
 \tag{T}
\]
Each is a first-order formula in the language with $I$; no further
inductive definition is introduced at an arrow. Quantification over
computable typed arguments abbreviates a finite disjunction over $S$.

We give the computability facts used in the constructor induction.
Joint induction on type proves that a non-lambda is computable whenever
all its direct reducts are computable; that variables and bottoms are
computable; that computability is preserved by $\leadsto$ and $R_S$;
and that computability implies $I$. At base type, closure of (A) and
transitivity of positive paths give the assertions. At an arrow, $QU$
has no top beta step when $Q$ is not a lambda. Apart from bottom its
reducts reduce the function $Q$, so codomain induction applies. Bottom
and variable applications are handled by this same neutral argument.
Function-side closure gives downward closure of (T).

For the implication from computability to $I$ at an arrow, apply the
term to a fresh computable variable $x$. Every positive path
$Q R_S Q'$ gives $Qx R_S Q'x$. Induction on $I(Qx)$ therefore proves
$I(Q)$: all descendants $Q'$ inherit the corresponding strictly
smaller accessibility premise.

The lambda fact is also needed. Suppose every computable instance
$V[x:=U]$ is computable. Substituting a fresh variable first gives
$I(V)$. Induct on this accessibility while retaining the universal
instance premise. The direct reducts of $(\lambda x.V)U$ are its given
beta instance, a function-side body reduct, or bottom. For a body
reduct, substitution takes a raw step to a raw path or equality; the
only possible loss of a step is harmless, as when a variable-to-bottom
step becomes bottom-to-bottom. Mark decreases commute with substitution
because their leading labels remain unchanged. Downward closure and the
induction hypothesis therefore make every reduct computable. The
neutral argument gives the application and hence the lambda.
All statements respect the finite equivalence quotient.

Next prove that accessibility of finitely many members implies $J$ of
their multiset. First prove closure of $J$ under multiset union by outer
$J$-induction on one component and inner $J$-induction on the other.
Label removed occurrences by their component of origin and assign each
new occurrence to a removed dominator. The target splits into two
components, each unchanged or decreased, with at least one decreased.
The nested induction and downward closure supply its $J$-membership.
Now $I$-induction on $Q$ gives $J(\{Q\})$: any replacement consists of
finitely many descendants $Q_i$, the induction hypothesis gives each
$J(\{Q_i\})$, and finite union gives the replacement's membership.
Another finite union proves the assertion for arbitrary accessible
member multisets.

Finally prove that a constructor with computable arguments is computable
by the lexicographic induction
\[
 (\text{head class},\ \text{base multiset},\
    \text{label in }\N\cup\{\omega\},\ \text{extra multiset}).
 \tag{L}
\]
Use finite induction on head classes, $J$-induction on the base list,
ordinary induction on finite labels followed by the $\omega$ case,
and $J$-induction on extras. Every earlier invariant quantifies all
later parameters, including the output type and the finite typed extra
list. These are first-order formulas in the language of $\ID$.

The neutral lemma reduces this assertion to its direct reducts.
Ordinary marking decreases $\omega$ to a finite label, and the added
lowering rule strictly decreases a finite label. Select applies a
computable selected argument to computable lower-label self-copies.
Copy first obtains those copies and then uses a smaller head class;
its target label may reset to $\omega$ because the first component
has decreased. In the multiset rule, every inserted argument is a
positive mark-decrease descendant of a removed one. Downward closure
makes it computable, and the base multiset decreases under $\DM(R_S)$.
This earlier decrease again permits the outer label to reset to
$\omega$. Abstraction creation decreases the label before adding a
fresh extra, and the lambda fact applies to every computable
substituted extra. A base or extra argument reduction decreases its
corresponding multiset. Bottom is computable. These exhaust the rules.

Structural induction under a computable substitution now makes every
auxiliary term computable and hence accessible. All syntax, typing,
paths, multiset assignments and invariants used above are arithmetic
in the one fixed universe, so this is an $\ID$ proof of the stated
individual memberships.
\end{proof}

\subsection{Finite budgets and preservation of the universe}
A symbolic marked term is $S$-admissible if its node and binder types
and every head's original declaration lie in $S$. Ordinary endpoints
are admissible for $S=\Ty(s)\cup\Ty(t)$.

\begin{lemma}[Locality]\label{lem:window-locality}
Every term in a finite weak or strict derivation with $S$-admissible
endpoints, and every numerical intermediate constructed below, is
$S$-admissible.
\end{lemma}
\begin{proof}
Induct over the actual comparison and occurrence-list derivations.
Var and matching introduce no type. Abs opens at the given binder
type. Put and Strict change only the mark. FAbs appends the domain
of the existing outer arrow and uses its codomain.

For Select, each supplied domain and every successive arrow suffix is
a subtype of the selected argument's type. A self-copy retains its
original head declaration. The full application consumes original
lambdas, substitutes admissible copies, and possibly marks a residual
head and appends those copies. Induction on that literal operation
introduces no new type or head. This establishes admissibility of the
continuation source before its derivation induction is applied.

Copy and Stat retype only to actual target-argument types. Their new
target head is present in the endpoint and therefore satisfies the
head-window guard. Their occurrence witnesses use actual endpoint
arguments. This proves the symbolic assertion.

Numerical realization retains this typed shape. Lowering and mark
decrease change no declaration. Copy and multiset reset use the
already admissible target head. Auxiliary Select applications and the
lambdas created by the full-tuple helper have the domain and suffix
types already checked above. Beta substitution, contexts and equivalent
head permutations preserve the same window. Thus every constructed
raw edge belongs to this single $R_S$ universe.
\end{proof}

Write $q\sim_n Q$ when $Q$ realizes the symbolic source $q$ with the
same typed shape, ordinary heads at label $\omega$, and every marked
head at some finite label at least $n$. The realization preserves base
and extra lists, commutes with simultaneous substitution, and weakens
to every smaller threshold. For an ordinary $q$ there is a unique
$\tr(q)$, available at every threshold.

\begin{lemma}[Full-tuple simulation]\label{lem:tuple-simulation}
If $\ap q{a_1,\ldots,a_l}$ is defined, sufficiently high realizations
of $q,a_i$ reduce by a finite auxiliary path to a realization of that
exact result at any prescribed lower threshold $n$. The helper spends
at most $l+1$ units beyond that threshold.
\end{lemma}
\begin{proof}
Consume the same original lambda spine by beta steps, with the supplied
arguments substituted in their common environment. If the tuple ends,
this is the required result. Otherwise source definedness says that
the original residual head is functional. Mark an ordinary head to
a sufficiently large finite label, or retain its available finite
label. For each remaining argument, create an abstraction and
beta-contract it, thereby appending that argument while decreasing
the label by one. There are at most $l$ such decreases. Substituted
old arguments retain their own thresholds. A protected-variable
overapplication is undefined in the source and is never a premise.
All helper types remain in the window by Lemma~\ref{lem:window-locality}.
\end{proof}

\begin{lemma}[Positive simulation]\label{lem:positive-simulation}
For every finite weak derivation $\W q t$ in a fixed window there is
a computable finite budget $b$ such that every $q\sim_b Q$ reduces,
possibly by an empty path and modulo finite equivalence, to $\tr(t)$.
For a strict derivation the budget can be chosen so that the path is
positive. Consequently, for ordinary confined endpoints,
\[
 s\cs t\quad\Longrightarrow\quad
 \tr(s)R_{\Ty(s)}\tr(t).
\]
\end{lemma}
\begin{proof}
All original targets are ordinary by Lemma~\ref{lem:elementary}.
Induct jointly on the finite comparison and list witnesses, taking
maxima of their finitely many required budgets.

Var requires no step. Abs lifts a body path; Fun performs child paths
inside base arguments and uses the finite head/permutation equivalence.
Put chooses a finite leading label above its child's budget.
Select reserves one unit for the Select rule, the full-tuple cost
of Lemma~\ref{lem:tuple-simulation}, and the continuation budget.
FAbs reserves one unit and lifts its body path. Copy reserves one
unit, builds its ordinary target from lower-label self-copies, and
performs the finite child paths in its arguments.

For Stat retain the specified source occurrences unchanged in the
intermediate target. Each inserted occurrence has a strict premise,
hence a defined mark and a shorter weak derivation from that mark.
If its numerical source argument is ordinary, choose an
$\omega$-to-finite mark above the child budget. If it is already
numbered, reserve an additional unit and strictly lower the finite
label. Under initial lambdas use the same context step. These are
exactly the permitted primitive replacements. The numerical multiset
rule builds an ordinary target whose retained arguments then follow
their weak paths and whose inserted arguments follow their marked
weak paths.

For Strict itself perform this same nonempty mark decrease on its
source before following its weak premise. At a symbolically
already-marked source the symbol stays the same but its finite
numerical label strictly decreases. This makes the simulated path positive. Label zero
is never reused because the budget reserves the needed unit.

All choices depend only on the finite proof certificate. Locality
keeps every intermediate in the original endpoint window. For
ordinary $s$, $\tr(s)$ has only $\omega$ labels and meets every
finite budget. Confinement makes that window precisely $\Ty(s)$.
The marked-target auxiliary relation of Appendix~\ref{app:composition}
is not part of this simulation; its reflexive marked strict judgments
are never interpreted as positive paths.
\end{proof}

We use the following published ordinal-analysis input. For a positive
arithmetic inductive operator, every individual membership atom provable
in $\ID$ has inductive norm strictly below the Bachmann--Howard ordinal.
With stages beginning at the operator applied to the empty set, the
norm for the accessibility operator is exactly descending rank
\cite[Definitions~6.3.1, 6.3.5 and 9.1.2,
Lemma~6.5.2, Theorems~9.5.5--9.5.6]{pohlers}.

\begin{theorem}[Well-foundedness and individual bound]\label{thm:upper}
The confined strict relation is well founded, and
$\rk_{\cs}(s)<\BH$ for every ordinary term $s$.
\end{theorem}
\begin{proof}
Every descendant of $s$ remains in $S=\Ty(s)$.
Lemmas~\ref{lem:arithmetic-accessibility} and
\ref{lem:positive-simulation} exclude an infinite strict chain.
Their positive-edge simulation gives
\[
 \rk_{\cs}(s)\le \rk_{R_S}(\tr(s)).
\]
The arithmetic accessibility proof supplies an $\ID$ proof of the
individual $I$-membership of $\tr(s)$ in this fixed window. The tagged
$I$ component has its standalone accessibility norm. Applying the
published norm theorem yields
$\rk_{R_S}(\tr(s))<\BH$, as required. No transitivity assumption
about the symbolic comparison is needed in this argument.
\end{proof}

\section{Native terms with cofinal lower ranks}\label{sec:lower}

For a closed native term $t:A$, let $r_A(t)$ be its rank in the
restriction of $\cs$ to closed native terms of that same type.
These ranks are countable by well-foundedness and countability of
syntax. Restriction cannot increase rank. The lower construction
uses these ranks themselves, without a rank borrowed from the
auxiliary numerical system.

\subsection{A carrier inside literal T}
Fix closed native terms $P,K:B\to B$ and $U:B$ with $\W K P$
and $\Ty(P)\subseteq\Ty(K)$. Put
$H=o\to B\to B$, $F=\lambda n:o.K$ and
\[
 W=R^B0FU,\qquad U_1=R^B0,\quad U_2=U_1F.
\]
For any required result type $\rho$ define proof intermediates
\[
\begin{aligned}
 Q_\rho&=@^\star_\rho(U_2,U),\\
 C_{2,\rho}&=@^\star_\rho(U_1,F;Q_B),\\
 C_{1,\rho}&=@^\star_\rho(R^B,0;C_{2,H},C_{2,B}),\\
 C_{0,\rho}&=(R^B)^\star_\rho(;C_{1,o},C_{1,H},C_{1,B}).
\end{aligned}
\]
They are marked comparison terms, not additional native constructors.

\begin{lemma}[Native carrier]\label{lem:carrier}
$W\cs U$, and $W\cs T$ implies $W\cs PT$ for every ordinary
closed $T:B$ carrying the stated support guard.
\end{lemma}
\begin{proof}
Select the second base argument of $Q_B=\mk(W)$ to obtain $U$.
For the second assertion, the hypothesis is a weak proof from $Q_B$
to $T$. Select the extras $C_{1,B},C_{2,B},Q_B$ in turn to prefix
that proof by a derivation from $C_{0,B}$.

To derive the other Copy premise, start from $C_{0,B\to B}$.
Select its extra $C_{1,H}$ with one self-copy of type $o$; this
appends that copy to the marked $C_1$ owner. Select its original
$C_{2,H}$ extra with one further self-copy, and then its original
base argument $F=\lambda n.K$ with one self-copy. The result is
literally $K$, since $n$ is unused. Append the given weak proof
from $K$ to $P$. Copy at $R^B>@$ now constructs $PT$ from these
two whole-source premises.

Three outer Select prefixes carry this proof back to $Q_B$.
They select $U_2,U_1,R^B$ with respectively one, two and three
simultaneous copies at their displayed types. Strict marking gives
$W\cs PT$. The support of $T$ is contained in that of $W$ by
hypothesis; the support of $P$ occurs through $K$, and the final
application adds only $B,B\to B$, already present. Thus every
comparison is confined. No composition theorem is used.
\end{proof}

Let the semantic base type be $\iota$, interpreted as the countable
ordinals, and translate it to the actual type $\sigma=o\to o$.
Translate arrows homomorphically. Define erasable templates
\[
 E_\iota(i)=i,\qquad E_{A\to D}(i)=\lambda x:|A|.E_D(i).
\]
For $B=|A|$, $\eta=B\to B$, and fixed $P:\eta,U:B$, put
\[
 K_i=((\lambda a:\eta.\lambda b:\eta.a)P)E_{A\to A}(i),
 \qquad
 W_i=R^B0(\lambda n:o.K_i)U.
\]

\begin{lemma}[Indexed carrier]\label{lem:indexed-carrier}
$W_i\cs U$, and $i\cs j$ implies both
$W_i\cs W_j$ and $W_i\cs PW_j$.
\end{lemma}
\begin{proof}
First $K_i\cs P$ has a direct proof. Mark the outer application and
select its first base argument with one self-copy. The result is the
marked inner application retaining $P$ as its second base argument.
Select $P$ with an empty tuple. Since $P$ is a subterm, confinement
also holds. Strict inclusion supplies the weak hypothesis of
Lemma~\ref{lem:carrier}.

The context defining $W_i$ contains $i$ in its erasable argument.
Context stability therefore gives $W_i\cs W_j$ from $i\cs j$.
Apply the carrier lemma to that target to get $W_i\cs PW_j$.
The common context contributes the same types; the support of $j$
is contained in that of $i$. The final application uses only types
already present. The selector proof did not compose beta reductions.
\end{proof}

\subsection{Rank realization and full-prefix iteration}
Semantic arrow types contain all set functions of the indicated types,
not only monotone ones. Countable suprema are pointwise at arrows.
For closed actual terms define
\[
 t\models_\iota\alpha\iff r_\sigma(t)\ge\alpha,
\]
and let $t\models_{A\to D}f$ mean that for every semantic $x:A$
and every closed $u\models_A x$, the ordinary application $tu$
realizes $f(x)$.

\begin{lemma}[Realization closure]\label{lem:realization}
Realization is preserved when a term is replaced by a strict predecessor.
It is downward closed in its semantic value, with pointwise order at
arrows. If one fixed term realizes every member of a nonempty countable
family, it realizes their pointwise supremum. The typed lambda
fundamental lemma holds for these realizers.
\end{lemma}
\begin{proof}
Induct on the semantic type. At $\iota$, use the rank inequality
for a strict edge and the defining property of ordinal supremum.
At an arrow, apply the induction hypotheses to each realizing
argument; context stability gives $tv\cs uv$ from $t\cs u$.
The other assertions are pointwise. For the lambda lemma, its
application beta-reduces strictly to the substituted body, so strict
predecessor closure transports the body's realization backward.
An induction on the number of beta steps permits any finite sequence.
\end{proof}

The closed native iterator is
\[
 \It_A=\lambda i:\sigma.\lambda p:\eta.\lambda u:B.\
 R^B0\bigl(\lambda n:o.((\lambda a:\eta.\lambda b:\eta.a)p)
                   E_{A\to A}(i)\bigr)u.
\]
For an arbitrary semantic $f:A\to A$ define
\[
 f^0(x)=x,\qquad
 f^{\beta+1}(x)=f(f^\beta(x)),\qquad
 f^\lambda(x)=\sup_{\xi<\lambda} f^\xi(x)
 \quad(\lambda>0\text{ limit}).
 \tag{I}
\]

\begin{theorem}[Native finite-type iteration]\label{thm:iteration}
$\It_A$ realizes $(\alpha,f,x)\mapsto f^\alpha(x)$ with convention
\textup{(I)}, even when $f$ is not monotone.
\end{theorem}
\begin{proof}
Fix closed realizers $P\models f$ and $U\models x$. For each closed
index $i:\sigma$, prove that the single term $W_i$ realizes every
orbit value $x_\beta=f^\beta(x)$ for $\beta\le r_\sigma(i)$.
Use outer well-founded induction on $r_\sigma(i)$ and inner induction
on $\beta$.

At zero, $W_i\cs U$ supplies the conclusion. If
$\beta=\gamma+1\le r_\sigma(i)$, the rank equation gives a closed
native $j$ with $i\cs j$ and $r_\sigma(j)\ge\gamma$.
Otherwise all successor ranks would be at most $\gamma$, contrary
to $r_\sigma(i)\ge\gamma+1$. The outer induction makes $W_j$
realize $x_\gamma$. Hence $PW_j$ realizes $f(x_\gamma)$, and
$W_i\cs PW_j$ transports that realization to $W_i$.

At a nonzero limit the same $W_i$ realizes every earlier $x_\xi$.
There are countably many earlier stages. The supremum clause of
Lemma~\ref{lem:realization} gives $x_\beta$. This uses the entire
prefix; a cofinal subsequence would not suffice for an arbitrary
nonmonotone orbit.

If $i\models\alpha$ then $\alpha\le r_\sigma(i)$.
Finally transport the realization backward through the finitely
many beta steps from $\It_A iPU$ to this carrier. All ranks and
edges are those of the closed native confined relation.
\end{proof}

\subsection{Explicit ordinal operations and published calibration}
Let $Z=\lambda z:o.z:\sigma$ and define
\[
 s=\lambda a:\sigma.((\lambda u:\sigma.\lambda v:\sigma.u)Z)a.
\]
$Z$ realizes zero. The body of $sa$ lies strictly above $a$ by
outer-argument Select, so beta transport makes $s$ realize right
successor. Put
\[
 \operatorname{Cap}(P,U)=R^B0(\lambda n:o.P)U.
\]
The same carrier is strictly above every finite $P^kU$, by induction
using Lemma~\ref{lem:carrier}. Consequently
$w=\operatorname{Cap}(s,Z)$ realizes $\omega$.

Writing $\It=\It_\iota$, define the following closed abbreviations:
\[
\begin{aligned}
 \mathrm{add}&=\lambda a\,b.\It\,b\,s\,a,\\
 \mathrm{mul}&=\lambda a\,b.\It\,b\,(\lambda x.\mathrm{add}\,x\,a)\,Z,\\
 \mathrm{one}&=sZ,\\
 \mathrm{pow}&=\lambda a.\It\,a\,(\lambda x.\mathrm{mul}\,x\,w)\,\mathrm{one},\\
 \mathrm{fix}&=\lambda f\,z.\It\,w\,f\,(s z).
\end{aligned}
\]
Their variables have the unique indicated ordinal-realizer types.
They realize, respectively, $\alpha+\beta$, $\alpha\beta$,
$1$, $\omega^\alpha$, and $f^\omega(\zeta+1)$.
The formulas specify the order of the noncommutative operations:
multiplication iterates addition on the right, and exponentiation
iterates right multiplication by $\omega$.
Let $\mathrm{Next}=\mathrm{fix}\,\mathrm{pow}$, realizing the
function that takes $\zeta$ to the least epsilon number above it.

Put $A_0=\iota$ and $A_{l+1}=A_l\to A_l$. The finite-type
extractors of Simmons \cite[Defs.~2.10--2.14]{simmons2004}
are
\[
 B_l h h_l\cdots h_1
   =\Fix\bigl(\lambda\alpha.\ h^\alpha h_l\cdots h_1\,0\bigr),
 \qquad B_l:A_{l+2}.
 \tag{B}
\]
They have actual closed native realizers. For $l=0$, the inner body
is $\It_\iota\,\alpha\,h\,Z$. For $l\ge1$, it is
\[
 (\It_{A_l}\,\alpha\,h\,h_l)h_{l-1}\cdots h_1Z.
 \tag{B$_T$}
\]
Prefix the corresponding typed lambdas and use $\mathrm{fix}$ for
$\Fix$. The terminal $Z$ is present in every case. The iterator
theorem and the lambda fundamental lemma verify the realization
of (B). All names here abbreviate finite expressions in the original
native constants, lambdas and applications.

We state the published calibration precisely. Let $\Omega$ be the first
uncountable ordinal, set $a_0=0$ and $a_{n+1}=\Omega^{a_n}$, and let
$\Lambda=\sup_n a_n=\varepsilon_{\Omega+1}$.
Simmons's finite-type iteration uses exactly the full-prefix pointwise
convention (I), on the full function spaces. His helpfulness theorem
for every $B_l$ is supplied in
\cite[Theorem~5.10]{simmons2008}, completing the hypotheses used
in \cite[Theorem~6.1 and Corollary~6.2]{simmons2004}.
For the Veblen base pair put
\[
 N_B=B_1B_0\mathrm{Next},\qquad \epsilon_B=\Gamma_0=N_B(\omega).
\]
The displayed identity following
\cite[Lemma~6.6]{simmons2008} gives the last equality.
The resulting sequence
\[
 \Delta_0=\Gamma_0,\quad
 \Delta_1=N_B\Gamma_0,\quad
 \Delta_{l+2}=B_lB_{l-1}\cdots B_0N_B\Gamma_0
 \tag{D}
\]
satisfies $\Delta_n=\psi_B(a_n)$ by the cited calibration.
Here $\psi_B$ is the Veblen-hull collapse, not the arithmetic-hull
collapse attached to $(\mathrm{Next},\varepsilon_0)$.

For completeness, the hypothesis called tameness in that calibration
holds for the displayed towers: their recursive coefficient norms are
at most one, whereas every value in the critical base sequence is at
least $\Gamma_0>1$. This is the actual inequality needed in the proof;
no assertion that a critical ordinal is finite is used.

Replace the operations in (D) by their native realizers and call the
resulting closed terms $t_n$. Each has actual outer type $\sigma=o\to o$,
and
\[
 r_\sigma(t_n)\ge\psi_B(a_n).
 \tag{R}
\]
The native syntax alone supplies these witnesses.

\subsection{The endpoint is the same Bachmann--Howard ordinal}
We record why the tower values are cofinal at the endpoint used for the
upper bound. In the Veblen convention, an $N_B$-critical ordinal is an
orbit value $N_B^\alpha(\Gamma_0)$, not a fixed point of the strictly
inflationary function $N_B$. These criticals are closed under countable
increasing suprema. Section~5 of \cite{simmons2004} characterizes
$\psi_B(\xi)$ as the least such critical $\theta$ for which
\[
 \eta<\xi,\quad |\eta|<\theta
       \quad\Longrightarrow\quad \psi_B(\eta)<\theta .
 \tag{P}
\]
The coefficient norm is recursively the maximum of the countable
coefficients in the base-$\Omega$ representation and of the norms
of its exponents. In particular $|a_n|\le1$.
The Veblen-hull endpoint $\psi_B(\varepsilon_{\Omega+1})$ is the
$\BH$ of the inductive-norm calibration
\cite[Sections~9.4--9.5]{pohlers}.

\begin{lemma}[Tower cofinality]\label{lem:tower-cofinality}
$\sup_n\psi_B(a_n)=\psi_B(\Lambda)=\BH$.
\end{lemma}
\begin{proof}
Let $\theta=\sup_n\psi_B(a_n)$. For $n<m$, the norm bound
$|a_n|\le1<\psi_B(a_m)$ and (P) give
$\psi_B(a_n)<\psi_B(a_m)$. The countable supremum is therefore
an $N_B$-critical ordinal.

If $\eta<\Lambda$ and $|\eta|<\theta$, choose $n$ large enough
that both $\eta<a_n$ and $|\eta|<\psi_B(a_n)$. Such a common
index exists by the two increasing cofinal sequences. Applying
(P) at $a_n$ gives $\psi_B(\eta)<\psi_B(a_n)\le\theta$.
Thus $\theta$ satisfies the endpoint clause, and its least
critical solution obeys $\psi_B(\Lambda)\le\theta$.
Conversely (P) at $\Lambda$ applies to each $a_n$, since
$|a_n|\le1<\psi_B(\Lambda)$, and gives
$\psi_B(a_n)<\psi_B(\Lambda)$. Take their supremum.
This uses the defining clause of this particular collapse; it
does not assume monotonicity for some different collapse convention.
\end{proof}

\begin{theorem}[Exact descending height]\label{thm:height}
The descending height of $\cs$ is $\BH$, already cofinally witnessed
by closed native T terms of outer type $o\to o$.
\end{theorem}
\begin{proof}
The realizers (R) and Lemma~\ref{lem:tower-cofinality} give the
lower bound for the restriction to closed native terms. Its ranks
are bounded by the ranks in the full ordering.
Theorem~\ref{thm:upper} bounds each individual full rank below
$\BH$. Taking the supremum of ranks plus one gives equality.
No individual term has been assigned rank $\BH$.
\end{proof}


\section{Conclusion and scope}
The preceding results prove Theorem~\ref{thm:main}: structural stability
and composition are established in Theorem~\ref{thm:stable};
Theorem~\ref{thm:decide} decides the same pair;
Theorems~\ref{thm:rpo} and~\ref{thm:orient} give its two rewriting
interfaces; and Theorems~\ref{thm:upper} and~\ref{thm:height} establish
the matching ordinal bounds.

The type-support condition is local to each source. It permits arbitrary
finite types in the input language, while keeping every descendant of one
source inside one arithmetic accessibility universe. The native family
used for cofinality ranges over unbounded finite type instances, which is
consistent with every individual rank being strictly below $\BH$.

The exact first-order restriction concerns actual ranked constructors.
It is not an uncurrying theorem for native $D$ or $R$. The proof also fixes
multiset status throughout. Other status policies require their own
composition and ordinal analyses.

The numerical comparison problem for Wilken--Weiermann assignments has
different inputs and a different algorithm. No collapse-control choice,
residue-table theorem or arithmetic quantifier-elimination procedure is
used in the construction of $\cw,\cs$. Those results can therefore be
developed independently without changing this ordering.

\paragraph{Preparation and verification.}
This manuscript was prepared with AI assistance and presents hand proofs
for the explicitly defined persistent-mark relation. A complete formal
verification of this relation and its ordinal calibration is not claimed.
The motivating marked-source obstruction concerns a particular earlier
formalized presentation; it is not a claim that Kop's published theorem
is false.

\appendix
\section{Composition through marked intermediate terms}
\label{app:composition}

This appendix proves composition for the relation actually defined in
Section~\ref{sec:comparison}. Auxiliary comparisons may have marked
targets; they will be eliminated when the final target is ordinary.
They are not used as a well-founded relation. In particular, an
already-marked term is strictly comparable to itself in this auxiliary
relation, so its strict part must not be substituted for the numbered
relation of Section~\ref{sec:upper}.

\subsection{The auxiliary rules}

A \emph{prefix copy} of
$P_\tau=f^\star_\tau(\vec b;\vec e)$ is
$f^\star_\rho(\vec b;\vec e_{\le j})$, where
$0\le j\le|\vec e|$ and $\rho$ is its required output type.
The prefix is literal, including order and types. Write $P+\vec a$
for the same marked owner with $\vec a$ appended to its extras.

The least finite-derivation pair $\succeq_e,\succ_e$ has all eight
rules of Section~\ref{sec:comparison}, with the subscript $e$ in
every premise, and the same persistent strict rule. It also has
the following seven weak rules. The base matching in \textsc{In}
allows a permutation; its extra matching does not.
\begin{enumerate}
\item[\textsc{In}]
From $f\approx g$, a weak matching of the base occurrences, and
an ordered, equal-typed weak matching of the extras, infer
\[
 \EW{f^\star_\tau(\vec b;\vec e)}
    {g^\star_\tau(\vec c;\vec d)}.
\]
\item[\textsc{Marked1}]
If $f>g$ and $P$, retyped as necessary, weakly dominates every
base and extra argument of $g^\star_\tau(\vec c;\vec d)$,
infer $\EW{P_\tau}{g^\star_\tau(\vec c;\vec d)}$.
\item[\textsc{Marked2}]
Use the same conclusion and whole-source premises when
$f\approx g$ and
$\vec b>_{\mathrm{mul}}^{\succeq_e,\succ_e}\vec c$.
\item[\textsc{Drop}]
From $\EW P t$ infer $\EW{P+\vec a}t$, for any typed extra list.
\item[\textsc{Perm}]
The source extras may be permuted, with each type attached to
its term.
\item[\textsc{Copies}]
If every entry of $\vec p$ is a prefix copy of $P$, then
$\EW{P+\vec p}t$ implies $\EW P t$.
\item[\textsc{PSelect}]
The \textsc{Select} rule may use a tuple of prefix copies of its
owner in place of the tuple of whole-owner copies. The selected
term remains an actual base or extra argument, and its full
original-spine call must succeed.
\end{enumerate}
Every original derivation embeds in this pair. Structural
induction using \textsc{Var}, \textsc{Abs}, \textsc{Fun}, and
\textsc{In} proves $\EW q q$ for every marked or ordinary term.

\subsection{A syntactic refinement relation}

We give explicitly the syntax relation needed to eliminate the
extra rules. It is useful because a staged special application
can expose one of the copies substituted at an earlier stage.
This situation is a syntactic refinement of the owner even
when it is not one full application of the initially selected
term.

Define $q\sharpens s$ by the following mutually inductive clauses.
They are finite trees, not a greatest fixed point.
Variables refine only themselves; lambdas refine componentwise;
ordinary heads have exactly the same symbol and positionwise
refining base lists. At a marked head,
\[
 f^\star_\tau(\vec b;\vec e)
       \sharpens f^\star_\tau(\vec c;\vec d)
 \tag{S}
\]
holds if $\vec b\sharpens\vec c$ positionwise and every typed
occurrence $d$ in $\vec d$ has one of two finite witnesses:
\begin{enumerate}[label=(\roman*)]
\item an actual occurrence $e$ of the same type in $\vec e$
with $e\sharpens d$;
\item $f^\star_{\operatorname{type}(d)}(\vec b;\vec e)
\sharpens d$.
\end{enumerate}
The second alternative is a premise tree of smaller height in
this mutual definition. No order or multiplicity matching of
extras is required. Head equivalence is insufficient in (S):
the symbol itself is preserved.

\begin{lemma}[Refinement syntax]\label{lem:sharp-syntax}
The relation $\sharpens$ is reflexive and transitive. It commutes
with renaming, opening, and simultaneous substitution by
pointwise refining environments. It preserves retyping of a
marked root and reflects a successful partial mark:
\[
 q\sharpens s,\quad \mk(s)=s'
 \quad\Longrightarrow\quad
 \mk(q)=q'\text{ and }q'\sharpens s'
 \quad\text{for some }q'.
 \tag{M}
\]
Suppose $q\sharpens s$ and $a_i\sharpens b_i$ in corresponding
types. If $\ap s{\vec b}$ succeeds, then $\ap q{\vec a}$ succeeds
and its result refines the result on the right.
\end{lemma}
\begin{proof}
Reflexivity is induction on the term; a marked extra uses its
own occurrence and the induction hypothesis. Transitivity is
mutual induction on the right refinement tree, simultaneously
for bases and the two kinds of extra witness. In the existing
extra case, invert the matching extra of the first refinement.
If it was supplied by an existing occurrence, compose the two
smaller refinements. If it was supplied by the whole owner,
compose that smaller witness with the right extra witness.
In the whole-owner case, apply the marked induction hypothesis
at the requested retyping. The resulting extra witness again
has exactly one of the two forms (i)--(ii). The variable,
lambda, and ordinary-head cases just compose their components.
This proves transitivity without any comparison judgment.

For two-environment substitution, induct over the same mutual
syntax tree. A variable uses its environment witness; bound
variables use a common fresh variable. Ordinary heads and
marked bases retain their positions. An existing extra remains
an existing substituted occurrence. A whole-owner witness is
the corresponding smaller refinement after substitution; the
substituted owner is precisely the owner of the substituted
extra list. Renaming and opening are instances of this proof.
Retyping leaves every base and extra premise unchanged.

For (M), refinement has the same root constructors on both
sides. Descend through their common lambdas. An ordinary head
is changed into the same marked head on both sides, with empty
extras. Marked heads are unchanged by both marks. A variable
on the right cannot have a successful mark.

For full application, follow the common original lambda
spines. If the tuple stops within the spine, the two-environment
substitution result applies. If it proceeds to a head, success
on the right excludes a variable head. Corresponding ordinary
or marked heads therefore remain, and the unconsumed arguments
are appended on both sides. Their positionwise refinements
supply existing-extra witnesses. Old existing-extra witnesses
remain valid after appending. In an old whole-owner witness,
the enlarged source refines the old source by retaining only
the old extras; compose this with that witness. All substitutions
within old bases and extras have already been covered.
\end{proof}

Several useful instances deserve to be stated. A marked owner
refines itself with any extras omitted or permuted. It refines
itself with any finite list of its prefix copies appended:
the old extras use (i), and each appended prefix uses (ii)
followed by the preceding omission instance. Conversely, an
owner with more extras refines the original owner. These
statements allow arbitrary typed multiplicities.

\begin{lemma}[Two-stage prefix calls]\label{lem:prefix-escape}
Let $\vec a,\vec b$ be lists of prefix copies of one marked
owner $P$. Suppose
\[
 \ap q{\vec a}=r,\qquad \ap r{\vec b}=w.
\]
Then either $\ap q{\vec a,\vec b}=w$, or $P\sharpens w$,
at the output type of $w$.
\end{lemma}
\begin{proof}
Follow the original initial lambda spine of $q$. When the
combined tuple does not pass its end, ordinary simultaneous
substitution composition gives the asserted equality. When
it passes that end and the original residual head is
functional, all remaining arguments are appended to that
same head, again giving the equality.

The only remaining case is that the combined call exposes an
original variable. A free variable cannot account for two
successful stages: a still-unconsumed argument at such a
variable would block the second stage. A bound variable can
have received an argument from the first batch. Its value is
a prefix copy of $P$, hence a marked functional owner.
The second call can act on that value. Its result is the
chosen prefix copy, retyped, with a suffix of the later
prefix-copy arguments appended. Every retained old extra is
an existing extra of $P$; every appended copy is refined by a
whole $P$. Clause (S) therefore proves $P\sharpens w$.
This is precisely the case in which full and iterated unary
application would disagree.
\end{proof}

\subsection{Eliminating the auxiliary rules}

\begin{lemma}[Ordinary-target elimination]\label{lem:conservative}
If $q\sharpens s$, $\EW s t$, and $t$ is entirely ordinary,
then $\W q t$. The corresponding statements hold for the
strict judgment, weak matching, strict multiset replacement,
and replacement suppliers.
\end{lemma}
\begin{proof}
Use mutual induction over the five finite comparison and
list-judgment trees. Keep the refinement witness arbitrary.
In \textsc{Var} and \textsc{Abs}, inversion of refinement
gives the same variable or the two body witnesses.
\textsc{Fun} uses the component induction hypotheses.
\textsc{Put} uses (M). For \textsc{Copy}, retype the whole
source witness and apply the induction hypothesis to each
target argument; the head is unchanged. \textsc{Stat} also
uses the multiset induction hypothesis. Opening commutes
with refinement, so \textsc{FAbs} follows from its body
induction hypothesis and the same rule.

For either selection rule, let $a$ be the selected argument
of $s$. If it is a base, refinement supplies an actual
corresponding base $b$ of $q$ with $b\sharpens a$. If it is
an extra, its witness either supplies such an actual extra
$b$ or says that the retyped whole $q$ refines $a$.
Every whole $q$ copy refines the corresponding prefix copy
of $s$: use $q\sharpens s$, then omit the unwanted extras.
In the actual-argument case, Lemma~\ref{lem:sharp-syntax}
gives a successful full call on $b$ with whole $q$ copies
whose result refines the given selected result. Apply the
continuation induction hypothesis and the original
\textsc{Select} rule.
In the whole-owner case, the same application lemma gives
an owner with whole copies appended, refining the selected
result. The original owner refines this expanded owner by
the instances following Lemma~\ref{lem:sharp-syntax}.
Transitivity of refinement therefore lets the continuation
induction hypothesis finish directly.

In \textsc{Drop}, the enlarged owner refines the owner in
the premise by omitting its appended extras. In \textsc{Perm}
it refines the owner with the premise order. In \textsc{Copies},
the unexpanded owner refines the expanded one by the prefix
copy instance. Compose the supplied refinement with this
syntactic witness and use the smaller premise induction
hypothesis.
The conclusions of \textsc{In}, \textsc{Marked1}, and
\textsc{Marked2} are marked and so cannot occur at this
ordinary-target weak judgment.

For a strict judgment, (M) gives the successful mark on the
refining source, and its weak premise is smaller. In the
list judgments, keep the very same partitions and
occurrence maps, apply the weak or strict induction
hypothesis at each occurrence, and retain nonemptiness of
the removed list. These cases complete the mutual induction.
\end{proof}

\begin{lemma}[Compression of extras]\label{lem:extra-compression}
Let $P=f^\star(\vec b;\vec e)$ and
$T=g^\star(\vec c;\vec d)$ have the same output type.
Suppose $f\approx g$, the base lists have an auxiliary weak
matching, and $\EW{P_{\operatorname{type}(d)}}d$ for every
target extra $d$. Then $\EW P T$.
\end{lemma}
\begin{proof}
Let $\vec p$ have, in the order of $\vec d$, one whole
$P_{\operatorname{type}(d)}$ for each occurrence $d$.
\textsc{In} first compares
$f^\star(\vec b;\vec p)$ with $T$.
\textsc{Drop} appends $\vec e$ to this source;
\textsc{Perm} moves $\vec e$ in front of $\vec p$.
Every entry of $\vec p$ is a full-prefix copy of
$f^\star(\vec b;\vec e)$, so \textsc{Copies} discharges
them and gives the assertion. No comparison composition
is used.
\end{proof}

\begin{lemma}[Refinement admissibility]\label{lem:sharp-admissible}
For arbitrary targets,
\[
 q\sharpens s,\quad \EW s t\quad\Longrightarrow\quad\EW q t.
 \tag{A}
\]
The strict and three list-judgment analogues also hold.
\end{lemma}
\begin{proof}
Repeat the mutual derivation induction of
Lemma~\ref{lem:conservative}, returning auxiliary judgments.
All its original-rule, macro, selection, and strict
constructions remain valid with a marked final target.
There are now three additional root cases.

For \textsc{In}, the source refinement preserves the
head and refines each base, so the base matching follows
from the induction hypotheses. A source extra matched
with a target extra $d$ has either an actual-extra
refinement witness $e_q\sharpens e_s$, or a whole-$q$
witness. In the first case the smaller matching
induction hypothesis gives $\EW{e_q}d$; zero-argument
\textsc{Select} then gives $\EW q d$. In the second case
the smaller hypothesis already gives $\EW q d$.
Lemma~\ref{lem:extra-compression} reconstructs the full
marked target with precisely its original extras.

For \textsc{Marked1}, transport each whole-source
premise by the weak induction hypothesis, at its stated
type, and reapply the rule. For \textsc{Marked2}, do
the same and transport its strict multiset by the list
induction hypothesis. This exhausts the extra cases.
Persistence matters in the strict induction: a marked
source is returned by its mark, so (M) applies there
without requiring a fresh syntactic mark.
\end{proof}

\begin{lemma}[Marking transport]\label{lem:extended-marking}
For all auxiliary terms,
\[
 \EW s t\Longrightarrow\EW{\headmark(s)}{\headmark(t)},\qquad
 \EW{\headmark(s)}t\Longrightarrow\EW s t.
 \tag{T}
\]
If $\EW s t$ and $\mk(t)$ is defined, then $\mk(s)$ is defined.
\end{lemma}
\begin{proof}
For the second implication descend through a common
initial lambda spine. At an ordinary head use
\textsc{Put}; at a marked head or a variable spine
the total mark is unchanged. A comparison whose source
is a lambda can use only \textsc{Abs}; a comparison
whose source is a variable is its \textsc{Var}
comparison. This inversion justifies the descent and
also proves reflection of markability.

Prove the first implication by mutual derivation
induction. \textsc{Fun} becomes \textsc{In} with empty
extras. \textsc{Copy} and \textsc{Stat} become,
respectively, \textsc{Marked1} and \textsc{Marked2},
with the same premises. Already-marked target rules
retain their conclusions. \textsc{Put} uses
idempotence. For either selection rule, its smaller
continuation gives the marked selected result above
the marked target. The second implication of (T)
removes the mark from that selected result, so the
same selection applies. \textsc{Abs} and
\textsc{FAbs} use the body hypotheses; the extra-list
macros lift their premise hypotheses. The strict and
multiset cases use the same occurrence witnesses
and $(\headmark)^2=\headmark$.
\end{proof}

\subsection{Covariance of actual full applications}

\begin{lemma}[Two environments]\label{lem:two-environments}
If $\EW s t$ and type-preserving environments $\gamma,\delta$
satisfy $\EW{\gamma(x)}{\delta(x)}$ for every variable,
then $\EW{s\gamma}{t\delta}$. The strict and list-judgment
versions hold as well.
\end{lemma}
\begin{proof}
Induct simultaneously over those judgments, with the
environments arbitrary. A variable uses the supplied
comparison; a binder extends both environments by
the same fresh variable. Head rules preserve all
their finite occurrence maps and use the corresponding
induction hypotheses.

In a selection rule the successful source application,
after $\gamma$, is exactly the original application
result after $\gamma$. This is the one-way
successful-application substitution identity from
Definition~\ref{def:application}. Whole and prefix
copies become literal copies of the substituted owner.
The continuation induction hypothesis has target
$t\delta$, so the same selection rule applies.
\textsc{FAbs} uses lifted environments. The extra-list
macros retain their list and prefix witnesses under
$\gamma$. Finally, a successful mark commutes with
substitution, including its identity case on a marked
root. It supplies the strict induction step. Equal
substituted terms do not collapse indexed occurrences.
\end{proof}

\begin{lemma}[Application covariance]\label{lem:application-covariance}
Suppose $\EW q s$ and $\EW{a_i}{b_i}$ in the respective
argument types. If $\ap s{\vec b}$ succeeds, then
$\ap q{\vec a}$ succeeds and
\[
 \EW{\ap q{\vec a}}{\ap s{\vec b}}.
 \tag{C}
\]
\end{lemma}
\begin{proof}
First check the domains, before decomposing a full call.
Failure is possible precisely when the tuple passes
the original lambda spine and exposes an original
variable. If a source has this form, inversion of its
auxiliary comparison forces the identical lambda
spine and variable on the target: only \textsc{Abs}
and then \textsc{Var} are available. Thus right-hand
full-call success implies left-hand full-call success,
independently of the actual well-typed argument values.

We next prove the one-argument assertion by induction
on $\EW q s$. The \textsc{Var} case is impossible.
For \textsc{Abs}, apply
Lemma~\ref{lem:two-environments} to the body with
$a,b$ assigned to the consumed variable.
\textsc{Fun} and \textsc{In} append $a,b$ to their
respective extra lists and use \textsc{In}.
\textsc{Put} has the same applied source as its
smaller marked-source premise.
\textsc{FAbs} substitutes $a,b$ for the fresh
variable in its premise, again by the two-environment
lemma.

In \textsc{Copy} or \textsc{Marked1}, the left
application appends $a$. The old whole-source
premises survive by \textsc{Drop}, and the new
target extra $b$ is dominated by selecting the new
extra $a$ and using $\EW a b$. \textsc{Marked1}
reconstructs the applied target. For \textsc{Stat}
or \textsc{Marked2}, use \textsc{Marked2} with
the unchanged base multiset. In \textsc{Drop}
and \textsc{Perm}, apply the induction hypothesis
and move the new argument past the old finite extra
list using \textsc{Perm}. In \textsc{Copies}, apply
the hypothesis to the expanded owner, move its
copies after the new argument, and discharge them:
each old prefix is still a prefix of the owner
with the new argument appended.

It remains to check selection. Let $P$ be its owner,
$v$ the selected argument, $\vec p$ its tuple of
prefix copies, and
$\ap v{\vec p}=r$ its result, with $\EW r s$.
Put $Q=P+a$, retyped as necessary. Selecting the
new extra gives $\EW Q b$. Apply the induction
hypothesis to the smaller continuation $\EW r s$,
using the whole $Q$ as left argument and $b$ on
the right. It gives
$\ap r Q=r'$ and $\EW{r'}{\ap s b}$.
Both $\vec p$ and the final whole $Q$ are prefix
copies of $Q$. Lemma~\ref{lem:prefix-escape} gives
either a successful full call $\ap v{\vec p,Q}=r'$
or $Q\sharpens r'$.
In the first case \textsc{PSelect} for the actual
argument $v$ of $Q$ gives the desired comparison.
In the second case
Lemma~\ref{lem:sharp-admissible} gives it directly.
This step uses no comparison cut.

Finally, a successful full call agrees with its
sequence of unary calls. Within the original
lambda spine this is ordinary substitution
composition; beyond a functional head it is
concatenation of extras. The blocked-variable
case has already been excluded for both full
calls. Repeatedly apply the unary result and
use these two identities to obtain (C).
\end{proof}

\subsection{Target transformations and the cut induction}

\begin{lemma}[Target extras]\label{lem:target-macros}
If $\EW s P$ for a marked target $P$, the conclusion
remains true when a suffix of its extras is removed,
when its extras are permuted, or when a finite list
of prefix copies of $P$ is appended.
\end{lemma}
\begin{proof}
For each assertion induct on the comparison derivation.
Only \textsc{In}, \textsc{Marked1}, and
\textsc{Marked2} create the marked target;
the other possible root rules retain their target
and lift the shorter induction conclusion.

For suffix removal in \textsc{In}, remove the
matching source suffix, apply \textsc{In}, and
restore that source suffix by \textsc{Drop}.
For a permutation, permute the matching source
extras, apply \textsc{In}, and restore the source
order by \textsc{Perm}. In either marked precedence
rule, keep just the needed target-extra premises
or reorder their indexing.

For copy addition in \textsc{In}, append to the
source the corresponding source-prefix copies.
Each such source prefix compares with the
requested target prefix by \textsc{In}, using
the restricted old matching. The complete
expanded lists therefore compare by
\textsc{In}; \textsc{Copies} discharges the
added source copies.
For \textsc{Marked1} or \textsc{Marked2}, the
same root rule, with a truncated target extra
list, compares the whole source with each
requested target prefix. These are exactly the
new whole-source premises needed to append
the target copies. No transitivity is invoked.
\end{proof}

\begin{theorem}[Ordinary-final cut]\label{thm:cut}
If $\EW s t$, $\EW t u$, and $u$ is entirely ordinary,
then $\W s u$. The source and intermediate term
may contain arbitrary marks.
\end{theorem}
\begin{proof}
We first derive $\EW s u$. Induct lexicographically
on the size of the right derivation and then that
of the left derivation. Sizes include all five
mutual judgments and every whole-source premise.
Thus a weak premise contained in a right strict
or multiset witness has smaller first coordinate.
When that coordinate decreases, the newly
constructed left proof may have any finite size.

If the left root is \textsc{Put}, either selection
rule, or an extra-list macro, apply the induction
hypothesis to its smaller continuation with the
right derivation fixed, and lift the result by
the same rule. After these cases have been
removed, inspect the right root.

\paragraph{Var, Abs, and Fun.}
For \textsc{Var}, $t=u$, so keep the left proof.
For \textsc{Abs}, the remaining left root is
\textsc{Abs} or \textsc{FAbs}. Cut their opened
body premises using the smaller right body
derivation and reapply that root.
For \textsc{Fun}, a remaining left
\textsc{Fun} composes the two occurrence
matchings by smaller-right cuts. A left
\textsc{Copy} keeps its strict head comparison,
which respects head equivalence, and transports
each whole-source premise by those cuts.
A left \textsc{Stat} transports both its strict
multiset and its whole-source premises through
the right weak matching. The precise multiset
transport is given below.

\paragraph{Put.}
The marking lemma gives
$\EW{\headmark(s)}{\headmark(t)}$.
The right premise compares $\headmark(t)$ with $u$
and is smaller than its \textsc{Put} conclusion.
Cut, and then remove the left total mark by
Lemma~\ref{lem:extended-marking}.

\paragraph{Copy and Stat.}
The intermediate term is marked, so the remaining
left root is \textsc{In}, \textsc{Marked1},
or \textsc{Marked2}. Each such comparison may
be retyped at the root without altering its
premises. Cut it against each smaller right
whole-source premise. Transitivity of head
precedence, and its compatibility with head
equivalence, give \textsc{Copy} when the right
rule is \textsc{Copy}.

If the right rule is \textsc{Stat}, a left
\textsc{Marked1} likewise yields \textsc{Copy}.
A left \textsc{In} transports its weak base
matching into the right strict multiset.
A left \textsc{Marked2} composes its strict
multiset with the right one. The whole-source
premises have already been transported;
\textsc{Stat} then gives the conclusion.

\paragraph{FAbs and target macros.}
A remaining \textsc{In}, \textsc{Marked1}, or
\textsc{Marked2} comparison survives opening
both owners with the same fresh extra variable.
In \textsc{In} append its reflexive comparison.
In the two precedence rules keep old
whole-source premises by \textsc{Drop} and
select the fresh extra for its own premise.
Cut with the smaller right body derivation,
then use \textsc{FAbs}.
For a right \textsc{Drop}, \textsc{Perm},
or \textsc{Copies}, transform the target of
the left comparison by
Lemma~\ref{lem:target-macros} and cut with
the smaller right premise.
The right root cannot be \textsc{In},
\textsc{Marked1}, or \textsc{Marked2},
because $u$ is ordinary.

\paragraph{Transport of strict multiset occurrences.}
Each final target occurrence is an ordinary
subterm of $u$. Consider an element path
$\EW a b$ followed by $\ES b c$.
Write $\mk(b)=b'$. Markability reflection
supplies $\mk(a)=a'$, and marking transport
gives $\EW{a'}{b'}$. Its cut against the
smaller right weak witness $\EW{b'}c$ gives
$\ES a c$. For a strict-then-weak path,
retain the original successful mark of $a$
and cut its weak witness. For a
strict-then-strict path, marking transports
$\EW{\mk(a)}b$ to
$\EW{\mk(a)}{\mk(b)}$, since the first
marked term is fixed by total marking.
The smaller right strict witness then
gives the composite strict edge.

To compose the occurrence partitions, keep
exactly those originally kept occurrences
whose intermediate images are kept again.
Remove every other source occurrence.
This removed set contains the first
removed set and is nonempty when the
first multiset is strict. Each target
occurrence is traced backwards through
the two supplier maps. A path that is
not the unique path between kept
occurrences has a strict edge, so the
element transport just proved supplies
its strict witness. If only the second
factor is strict, its nonempty removed
set has a nonempty inverse image under
the first weak bijection. The
strict-then-weak case is the same
construction with the second removed
set empty. Unused removed occurrences
are allowed. This verifies all strict,
weak--strict, and strict--weak cases,
including deletion without insertion.

\paragraph{The selection continuation.}
Let the right rule select an actual
argument $q$ of a marked owner $Q$,
apply it to prefix copies $Q_j$, obtain
$r$, and continue by $\EW r u$.
This continuation has smaller right
size. Let $P$ be the source of the
remaining left comparison with $Q$.

If that comparison is \textsc{In},
the selected base or extra has an
actual matching occurrence $p$ of $P$
with $\EW p q$. For each target prefix
$Q_j$, \textsc{In} compares the
corresponding source prefix with it,
and \textsc{Drop} restores the
remaining source extras. A whole $P$
copy at the required type therefore
dominates $Q_j$.
Lemma~\ref{lem:application-covariance}
gives a successful full call on $p$
with whole $P$ copies and compares
its result with $r$. Cut with the
smaller continuation $\EW r u$,
then use the original \textsc{Select}
rule at the actual argument $p$.

If the left comparison is
\textsc{Marked1} or \textsc{Marked2},
its whole-source premises give
$\EW{P_{\operatorname{type}(q)}}q$.
Here the type is the \emph{full
function type} of $q$.
The same root clause with truncated
target extras compares a whole $P$
with every requested $Q_j$.
Application covariance now gives
\[
 \EW{\ap{P_{\operatorname{type}(q)}}
                 {\text{whole }P\text{ copies}}}{r}.
\]
For an empty tuple the left result is
$P$ itself; otherwise it is exactly
$P$ with those copies appended.
Cut with $\EW r u$ and discharge
the appended full-prefix copies by
\textsc{Copies}. This finishes the
selection case.

Every root has been covered, and
every invocation used the stated
lexicographic decrease. We have
obtained $\EW s u$.
Apply Lemma~\ref{lem:conservative}
with the reflexive refinement of $s$
to obtain the required original
judgment $\W s u$.
\end{proof}

\bibliographystyle{plain}
\bibliography{references}
\end{document}
